\documentclass[12pt,a4paper,oneside]{article}
\usepackage[brazil]{babel}
\usepackage[utf8]{inputenc}
\usepackage{olymp}
\usepackage{color}
\usepackage[dvipsnames]{xcolor}
\usepackage{listings}

\setcounter{problem}{2}

\begin{document}

\begin{problem}{Ano Bissexto}{bissexto}{2}

Faça um programa que, dado um certo ano, determine se ele é bissexto ou não.

O ano bissexto ocorre a cada quatro anos, nos anos múltiplos de 4 (exceto anos múltiplos de 100, que não são bissextos, a não ser os múltiplos de 400, que são). Veja alguns exemplos:

\vspace{-6pt}
\begin{itemize}
\item 2023 não é bissexto, pois não é múltiplo de 4
\vspace{-6pt}
\item 2020 foi bissexto e 2024 será também
\vspace{-6pt}
\item 2100 não será bissexto pois, apesar de múltiplo de 4, é também múltiplo de 100
\vspace{-6pt}
\item 2000 foi bissexto pois era múltiplo de 400, que sobrepõe a regra de não ser múltiplo de 100
\end{itemize}
\vspace{-6pt}

Para a resolução do problema construa \underline{obrigatoriamente} uma função com o protótipo abaixo, que recebe um inteiro representando o ano e retorna um valor lógico, sendo verdadeiro se o ano passado por parâmetro for bissexto, ou falso caso contrário:

\texttt{\textcolor{Mulberry}{bool} bissexto(\textcolor{Mulberry}{int} n)}

%\Notes
%
%Observações, se tiver.

\InputFile

A entrada contém várias linhas, cada uma contendo um ano de 4 dígitos compreendido entre 1582 e 3999.
A entrada termina \underline{com um ano negativo}, que não deve ser processado.

\OutputFile

Seu programa deve gerar uma linha de saída para cada ano informado na entrada, com a palavra ``\texttt{bissexto}'', caso o ano informado seja bissexto, ou a mensagem ``\texttt{nao bissexto}'', caso contrário.

% \Notes
% Nesta questão, seu programa deve usar a função \texttt{main} dada acima exatamente como está, não a modifique! Sua tarefa é apenas criar a função \texttt{eh\_primo}.

\Example

\begin{example}
\exmp{
1896
1900
1904
2000
2015
2016
2096
2100
2104
-1
}{
bissexto
nao bissexto
bissexto
bissexto
nao bissexto
bissexto
bissexto
nao bissexto
bissexto
}%
\end{example}


%Comentários sobre o exemplo, se necessário.

\end{problem}

\end{document}
